%% ma: L11-B16-0051
%% lop: 11
%% chuong: 5
%% bai: 16
%% dang: 11.16.2
%% loai: TN4
%% muc: VDC
%% dap_an: "D"
%% nguon: "tai-lieu-hoc-tap-mon-toan-lop-11-theo-chuong-trinh-giao-duc-pho-thong-2018.pdf (Lớp Toán Thầy Duy, Phạm Lê Duy, THPT Chu Văn An), Chương 3 Bài 2, Phần 1, Câu 39 (tr.372)"
%% kien_thuc: "Lớp 11 Bài 16; giới hạn lượng giác cơ bản lim sin x/x (ngoài SGK KNTT Bài 16, chương trình nâng cao)"
%% trang_thai: cho-duyet
%% ly_do: "Dạng mới đề xuất (11.16.x), chờ giáo viên duyệt dạng; tài liệu gốc không có lời giải mẫu, đáp án do kho tự giải và kiểm bằng sympy."
%% ghi_chu: "Câu dùng giới hạn (1−cos kx)/x^2 và sin x/x khi x→0, kiến thức không có trong SGK KNTT Bài 16; nhờ giáo viên quyết định có giữ lại."
\begin{ex}
Cho hàm số $y=f(x)=\dfrac{1-\cos3x\cos5x\cos7x}{\sin^27x}$. Tính $\lim\limits_{x\to0}f(x)$.
\choice{$\dfrac{83}{49}$}{$\dfrac{105}{49}$}{$\dfrac{15}{49}$}{\True $\dfrac{83}{98}$}
\loigiai{\sloppy Dùng $\lim\limits_{x\to0}\dfrac{1-\cos kx}{x^2}=\dfrac{k^2}2$ và $\lim\dfrac{\sin7x}{x}=7$. Phân tích $1-\cos3x\cos5x\cos7x=(1-\cos3x)+\cos3x(1-\cos5x)+\cos3x\cos5x(1-\cos7x)$, chia cho $x^2$ ta được giới hạn $\dfrac92+\dfrac{25}2+\dfrac{49}2=\dfrac{83}2$. Mẫu $\dfrac{\sin^27x}{x^2}\to49$. Vậy giới hạn $=\dfrac{83}{98}$.}
\end{ex}
